% V3.12 source-based contact and mount evidence. Source report and numerical
% results remain unchanged; figure data from archived accepted source states.
\clearpage
\section{Interstage copper, solder interfaces, and the TO mounting path (V3.12)}
\label{sec:v312-contact-mount}
The interstage metal paths are one of the largest existing construction sensitivities in this source-geometry cooler. At the accepted 0.400~A, 300~K reservoir, 100~Pa krypton, and 2.000~mW imposed device heat reference, the three formed-copper networks 58, 80, and stepped three-foot return~88 are conservative electrical/thermal multiports. Their total signed heat delivered to the hot, middle, and upper ceramics is \(-110.578\), \(+25.487\), and \(+87.173\)~mW, respectively. The sum is \(+2.082\)~mW, equal to the source Cu Joule generation; the 87.173~mW entering the upper ceramic is a \emph{redistributed} parasitic heat flow, not 87.173~mW of extra electric input.

\fig[.88]{v312_source_copper_heat_ledger}{pdf}{Original accepted source-geometry V3.5 native copper multiport heat balance. Negative signed values remove heat from the named ceramic. These three components cannot be added as if they were independent heat sources, and the values are not new fitted operating points.}

The upper stage receives \(+41.058\)~mW from net~80 and \(+46.116\)~mW from the upper foot of the stepped return net~88. In contrast, the four effective cavity gas branches transfer a total \(+3.420\)~mW from 100~Pa Kr and the cavity radiation approximation adds \(+0.563\)~mW, to \emph{several different thermal nodes}. The two copper supply leads exchange \(-11.873\) and \(+1.948\)~mW at separate hotter ceramic landings. These are source-specific heat-port values, not a directly additive thermal-load budget for the detector cold tip.

\subsection{Corrected BiSn foil-foot bondline and contact resistance}
Each formed foil has a CAD section \(0.300\times0.030\)~mm and each source joining foot has contact area \(A_c=0.060\ \mathrm{mm^2}\). Its existing effective thermal areal resistance is the \emph{assumed} \(r''_c=10^{-6}\ \mathrm{m^2 K/W}\). For an ideal, completely wetted single solder layer, the bulk-only limit is
\[
 r''_{\rm bulk}=\frac{t}{k_s},\qquad
 R_{\rm foot,bulk}=\frac{t}{k_s A_c}.
\]
The actual formed-foil source joint is \textbf{Bi57Sn43}, not the Sn95Sb5 used at the different P/N leg interfaces. Indium Corporation lists \(k_s=19\ \mathrm{W/(m K)}\) at 85\,$^\circ$C for the adjacent eutectic Bi58Sn42 (Indalloy~281), used here only as an \emph{alloy-family proxy at a different temperature}, and 61~W/(m K) for counterfactual SAC305. These are bulk-material reference values, \emph{not} qualified low-temperature interface data for this synthetic source. The nominal model resistance equals the ideal BiSn-family proxy bulk resistance of a \textbf{19~micrometre} bondline. By contrast, the earlier extreme \(10^{-4}\ \mathrm{m^2K/W}\) contact hypothesis is bulk-equivalent to \textbf{1.9~mm} of this solder and should not be mistaken for a realistic intact solder-foot tolerance. Real wetting, intermetallic layers, plated-film lateral spreading, and interface conditions remain unresolved. The SAC305 comparison is \emph{counterfactual}, not a claim about the source solder composition.

\begin{center}\small
\begin{tabular}{lrrrr}\toprule
Ideal-bulk solder hypothesis & Bondline (\(\mu\)m) & \(r''\) (\(10^{-6}\) m\(^2\)K/W) & Upper Cu (mW) & \(\Delta Q\) (mW)\\\midrule
BiSn proxy, \(k=19\) & 20 & 1.053 & 86.973 & -0.201\\
BiSn proxy, \(k=19\) & 25 & 1.316 & 85.985 & -1.189\\
BiSn proxy, \(k=19\) & 35 & 1.842 & 84.079 & -3.094\\
SAC305 (alternative), \(k=61\) & 20 & 0.328 & 89.825 & +2.651\\
SAC305 (alternative), \(k=61\) & 35 & 0.574 & 88.835 & +1.662\\\bottomrule
\end{tabular}
\end{center}
\fig[.82]{v312_solder_contact_fixed_field_screen}{pdf}{Source-native ceramic face temperatures are frozen while only the ideal bulk-only solder contact conductance varies. This is a conditional, fixed-field copper-flux screen, not a full-system detector-temperature sweep. All heat inputs are reconstructed from saved data without new FEM solves.}

Under these selected hypotheses, 20--35~micrometre BiSn-family proxy bonds change the source upper-stage Cu heat from 86.973 to 84.079~mW at fixed face temperatures, a 2.894~mW span. The original CAD joining foot is 20~micrometres thick; these larger values are a thickness experiment, not the as-built design. In comparison, the prior 100-fold hypothetical contact resistance produces a 6.276~K \emph{fully coupled} detector-temperature decrease; the latter is an extreme interface-isolation experiment, \emph{not} a manufacturable cooler improvement. Net~58 has source-segment thermal conductance 0.002087~W/K (about 479~K/W along the foil) versus 16.67~K/W per nominal foot, illustrating why ideal solder thickness alone is unlikely to eliminate the formed-Cu bypass.

\subsection{Matched-temperature TO mounting penalty and engineering priorities}
An independently accepted V3.8 finite-mount experiment, using its \emph{own provisional temperature-dependent detector Joule load}, increased external mounting resistance by 5 and 10~K/W \emph{in addition to} the existing 4.8334~K/W source package resistance. At fixed 0.400~A this raised the detector temperature from 218.900~K to 220.393 and 221.856~K. Bringing the detector back to within 0.032~K of the reference required currents of approximately 0.4153 and 0.4332~A and increased TEC terminal input from 0.595668~W to 0.648375 (+8.85\%) and 0.712481~W (+19.61\%). These mounting results must not be conflated with the fixed-2~mW solder screen.

The next physical qualification should therefore investigate \textbf{(1)} actual copper-foil, landed solder, and plated-film three-dimensional spreading with electrical continuity; \textbf{(2)} the installed TO mounting/contact geometry and thermal impedance; and \textbf{(3)} narrow exposed cavity-gas paths before geometry-consistent pellet-height and cross-section design optimization. The material and contact parameter intervals are test hypotheses rather than measured statistical tolerance distributions. Full source/audit hashes, the fixed-face solver, and original figures are available in the dedicated V3.12 technical record and its machine-readable result manifest.\footnote{Technical record: \texttt{docs/V3\_12\_COPPER\_JOINT\_MOUNTING\_EVIDENCE\_20261109.md}. Bulk properties: Indium Corporation, Solder Thermal Interface Materials (Bi58Sn42 Indalloy~281, 85\,$^\circ$C analog) and Indalloy~256 SAC305 preform data.}
